\documentclass[10pt,a4paper]{amsart}
\usepackage[margin=19mm]{geometry}
\usepackage{amsmath,amssymb,mathtools}
\usepackage[scr=rsfs,cal=euler]{mathalfa}
\usepackage{xfrac}
\usepackage[unicode,hidelinks,bookmarks=false]{hyperref}
\newcommand{\ie}{i.e.\ }
\newcommand{\cf}{cf.\ }
\newcommand{\resp}{respectively\ }
\newcommand{\Subsection}{\subsection}
\providecommand{\nobreakdash}{\nobreak\hbox{-}}
\setlength{\emergencystretch}{2em}
\multlinegap0pt
\usepackage{amssymb}
\usepackage{caption}
\usepackage{algorithm}\usepackage{algpseudocode}
\usepackage{mathtools}
\usepackage{multirow}
\usepackage{xcolor}
\usepackage{tikz}
\usepackage{csquotes}
\usepackage[normalem]{ulem}
\newtheorem{lemma}{Lemma}
\newcommand{\C}{\mathbb{C}}
\newcommand{\CH}{\operatorname{Ch}}
\newcommand{\PP}{\mathbb{P}}
\newcommand{\aA}{\mathcal{A}}
\DeclareMathOperator{\gr}{Gr}
\newcommand{\grk}{\operatorname{Gr}(k,k+2)}
\title{Positive Genus Pairs from Amplituhedra}
\author{Joris Koefler, Dmitrii Pavlov, Rainer Sinn}
\date{Source: arXiv:2601.11142v2 (CC BY 4.0)}
\begin{document}
\maketitle

\noindent\emph{Source and licence.} Positive Genus Pairs from Amplituhedra. Version arXiv:2601.11142v2 (CC BY 4.0), distributed by arXiv under the Creative Commons Attribution 4.0 International licence, http://creativecommons.org/licenses/by/4.0/, as recorded in the arXiv licence metadata for this exact version and shown on the abstract page https://arxiv.org/abs/2601.11142v2. Full licence notice: \texttt{licenses/arxiv-2601.11142-CC-BY-4.0.txt}.\par\smallskip

\noindent\textit{Editorial scope.} Counted unit: the article's lemma lem:m=4\_genus\_residual\_curve (source0.tex lines 791--799). The article's complete original proof of it is delivered with it (source0.tex lines 800--805, one \texttt{proof} environment of 6 lines). No mathematics is written, repaired or re-derived: the statement and the proof are the article's own lines, in the article's own notation, and no retained line is substituted or re-flowed.  Retained uncounted as the source-local context the proof needs: lines 190--193; lines 608--616; lines 636--638.  Nothing was omitted from any retained span, so the package carries no omission record.  No retained line contains a citation, so the wrapper carries no bibliography and no bibliography entry.  No retained line contains a figure environment, an image-inclusion command or drawing code, so no image is delivered and the package declares no asset dependency.  Editorial formatting only, in addition: an AMS \texttt{amsart} preamble that restates the article's own declarations and macros used by the retained lines, namely the environments \texttt{lemma} on separate counters, together with the standard packages those lines need.  The retained environments print in this document as Lemma 1 in order of appearance, whereas the article prints the same items with the numbering of its own full document.  The complete native source of the article as distributed in the arXiv e-print for this exact version is bundled unchanged as \texttt{sources/192835/source0.tex}.
\begin{align}\label{eq:Grkn_deg}
    \deg\gr(k,n)=(k(n-k))! \cdot 
        \frac{0! \cdot 1! \cdot \ldots \cdot (k-1)!}{(n-k)! \cdot (n-k+1)! \cdot \ldots \cdot (n-1)!}.
\end{align}

\begin{lemma}\label{lem:genus_residual_curve}
    %Let $D_i = \CH(L_i) \subset \gr(k,k+2)$ be the irreducible components of the algebraic boundary $\partial_a\aA_{n,k,2}$.
    Let $E$ be a curve obtained from the $(2k-1)$-fold transversal intersection of $D_i=\CH(L_i)$, where the lines $L_i\subset \PP(\C^{k+2})$ spanned by $Z_i$ and $Z_{i+1}$ are pairwise disjoint.
    Then, for $k\geq 3$ the (geometric) genus of $E$  can be computed as 
    \[
    g(E)=1+\frac{k-3}{2}C_k,
    \]
    where $C_k$ is the $k$-th Catalan number. For $k=1,2$ the genus of the curve $E$ is zero.
\end{lemma}

    \begin{align}\label{eq:can_div_curve}
    \omega_E\cong \mathcal{O}_{\grk}(-k-2)\otimes\mathcal{O}_{\grk}(2k-1)\cong \mathcal{O}_{\grk}(k-3).
    \end{align}

\begin{lemma}\label{lem:m=4_genus_residual_curve}
    Let $\mathcal{I}$ be a collection of $4k-1$ sets $I=\{i,i+1,j,j+1\}$ such that $I\cap I'=\varnothing$ for all $I,I'\in \mathcal{I}$.
    Then, let $E$ be a curve obtained from the $(4k-1)$-fold intersection $D_{\mathcal{I}}=\bigcap_{I\in \mathcal{I}}D_I$ of boundary divisors $D_{ij}$.
    Then, for $k\geq 2$ the (geometric) genus of $E$  can be computed as 
    \[
    g(E)=1+\frac{3k-5}{2}\deg(\gr(k,k+4)).
    \]
    For $k=1$ the genus of the curve $E$ is zero.
\end{lemma}

\begin{proof}
    The proof is identical to that of Lemma \ref{lem:genus_residual_curve} upon noting the following.
    Firstly, the canonical divisor of the curve $E$ obtained from $4k-1$ intersections of Schubert hyperplane sections, is given as $\omega_E\cong \mathcal{O}_{\gr(k,k+4}(-k-4)\otimes \mathcal{O}_{\gr(k,k+4)}(4k-1)\cong \mathcal{O}_{\gr(k,k+4)}(3k-5)$, as in \eqref{eq:can_div_curve}.
    Secondly, the degree of $\gr(k,k+4)$ now given by \eqref{eq:Grkn_deg} for $n=k+4$.
    Then, applying the Riemann--Roch theorem we obtain the result.
\end{proof}

\end{document}
